\section{Background}
\label{sec:background}

\subsection{Moving-Knife as a Procedural Metaphor}

Fair-division theory studies procedures for dividing a resource among
claimants. Dubins and Spanier's moving-knife protocol is one of the classic
continuous procedures~\citep{dubins1961}: a knife moves across the resource,
and the stopping rule determines the allocation.

BISECT's \mka{} does not import all of the envy-free or proportionality
guarantees from cake-cutting. Redistricting has fixed geography, indivisible
tracts, population balance constraints, contiguity requirements, and legal
criteria that do not map cleanly onto private valuations. The useful transfer
is procedural: a public, mechanical sweep is easier to explain and audit than
an opaque choice among many candidate maps.

\subsection{Reock Compactness}
\label{subsec:reock}

Reock compactness is
\begin{equation}
  \reock{}(D) = \frac{A(D)}{A(\mathrm{MEC}(D))},
  \label{eq:reock}
\end{equation}
where $\mathrm{MEC}(D)$ is the minimum enclosing circle of district
$D$~\citep{reock1961}. A compact round region fills much of its enclosing
circle; a long thin region does not.

The implemented scorer is an approximation. It computes the minimum enclosing
circle of projected tract centroids, not the exact district polygon boundary.
When explicit tract areas are not passed to the scorer, the numerator is the
number of centroid points. The resulting value is clamped to $[0,1]$.

\begin{remark}[Current score versus publication Reock]
\label{rem:centroid-approx}
The current split scorer is useful for ranking candidate sweep directions
inside the algorithm. It is not yet a final legal Reock statistic for a
certified district plan. A publication-grade Reock package should compute the
minimum enclosing circle of polygon boundary vertices and use true district
area in the numerator.
\end{remark}

\subsection{Minimum Enclosing Circles}
\label{subsec:welzl}

The implementation uses a deterministic-shuffle version of Welzl's randomized
incremental minimum-enclosing-circle algorithm~\citep{welzl1991}. For a set of
points in the plane, the enclosing circle is determined by at most three
boundary points. This makes the scoring routine fast enough to call twice per
candidate orientation.

\begin{center}
\begin{tabular}{c@{\quad}c@{\quad}c}
\textbf{Two-point basis} & \textbf{Three-point basis} & \textbf{Elongated set} \\
\begin{tabular}{ccc}
$\bullet$ &  & $\bullet$ \\
 & $\circ$ &
\end{tabular}
&
\begin{tabular}{ccc}
$\bullet$ &  & $\bullet$ \\
 & $\circ$ & \\
 & $\bullet$ &
\end{tabular}
&
\begin{tabular}{ccccc}
$\bullet$ & $\bullet$ & $\bullet$ & $\bullet$ & $\bullet$
\end{tabular}
\end{tabular}
\end{center}

The third sketch is the danger case \mka{} tries to avoid: a sweep that creates
an elongated half has a large enclosing circle and therefore a lower score.

\subsection{Where \mka{} Sits in the Portfolio}
\label{subsec:related}

\begin{table}[t]
\centering
\small
\caption{Structure-layer roles. ``Direct target'' means the implemented split
routine scores candidates by that quantity; it does not imply final-plan global
optimality.}
\label{tab:algorithm-compare}
\renewcommand{\arraystretch}{1.25}
\begin{tabular}{@{}p{0.16\linewidth}p{0.25\linewidth}p{0.21\linewidth}p{0.25\linewidth}@{}}
\toprule
Algorithm & Candidate mechanism & Direct target & Main caveat \\
\midrule
\metis{} & graph partition & edge cut & graph objective, not Reock \\
\cvdgeo{} & geographic assignment & seed proximity & half-split target today \\
\bfsg{} & two-frontier growth & legible contiguity & baseline, not compactness proof \\
\mka{} & projection sweep & centroid-MEC Reock & local and approximate \\
\bottomrule
\end{tabular}
\end{table}

\mka{} is therefore best understood as a diagnostic and construction option
for Reock-shaped geometry. It supplies a clear candidate ledger:
orientation, population prefix, left score, right score, selected minimum.
That ledger is the bridge to future audit artifacts.
